% !TeX program = xelatex
% 第123页第（2）题：全部不全等图形。
% 在安装了 ctex 的环境中使用 ctexart；最小 TeX 环境自动使用 article 兼容排版。
\newif\ifhasctex
\IfFileExists{ctexart.cls}{\hasctextrue}{\hasctexfalse}
\ifhasctex
  \documentclass[11pt,a4paper,UTF8,fontset=none]{ctexart}
\else
  \documentclass[11pt,a4paper]{article}
\fi
\usepackage{fontspec,amsmath,array,geometry,tikz,xcolor}
\geometry{left=12mm,right=12mm,top=12mm,bottom=13mm}
% Explicit mainland simplified-Chinese fonts, shared with the current source package.
\setmainfont{FandolKaiSC-Regular.ttf}[Path=fonts/,AutoFakeBold=2]
\ifhasctex
  \setCJKmainfont{FandolKaiSC-Regular.ttf}[Path=fonts/,AutoFakeBold=2]
  \newCJKfontfamily\ansfont{FandolKaiSC-Regular.ttf}[Path=fonts/]
  \newCJKfontfamily\headfont{NotoSansSC-Regular.ttf}[Path=fonts/]
\else
  \XeTeXlinebreaklocale "zh"
  \XeTeXlinebreakskip=0pt plus .1em minus .01em
  \newfontfamily\ansfont{FandolKaiSC-Regular.ttf}[Path=fonts/]
  \newfontfamily\headfont{NotoSansSC-Regular.ttf}[Path=fonts/]
\fi
\definecolor{redpen}{HTML}{C5001C}
\pagestyle{empty}
\setlength\parindent{0pt}
\setlength\parskip{0pt}
\newcommand\ans{\ansfont\color{redpen}}
\newcommand\pagetitle[2]{%
  {\headfont\fontsize{17pt}{22pt}\selectfont 第\(123\)页第（\(2\)）题：全部情况\par}
  \vspace{2mm}{\headfont\fontsize{12.5pt}{17pt}\selectfont #1\hfill 附页 #2\par}
  \vspace{2mm}\hrule\vspace{3mm}}
\newcommand\footer[1]{\vfill\hrule\vspace{1.5mm}{\fontsize{9pt}{12pt}\selectfont 小学数学同步练习册四年级上册\hfill 第\(123\)页附页 #1}}
\newcommand\pair[2]{\noindent #1\hfill #2\par}
\begin{document}
\pagetitle{一、梯形（共\(13\)种，本页为第\(1\)至第\(6\)种）}{1}
{\fontsize{10.5pt}{15pt}\selectfont
纸条长\(12\)厘米、宽\(2\)厘米，每小段长\(3\)厘米，上下各有\(5\)个交点。
以下红线均连接这些交点；每幅图中的灰色长方形都表示整张纸条。\par
\vspace{2mm}
\textbf{怎样算“不同”：}平移、旋转或翻转后能够\textbf{完全重合}的图形，只算一种。
梯形只有一组对边平行；为便于比较，统一把较短的底放在上方。\par}
\vspace{4mm}
\pair{\begin{minipage}[t][44mm][t]{90mm}\ans\fontsize{10.6pt}{14pt}\selectfont
{\headfont\fontsize{11pt}{14pt}\selectfont 梯形\(1\)}\par\vspace{1mm}
\begin{tikzpicture}[x=6.3mm,y=6.3mm,line join=round,line cap=round]
\path[use as bounding box] (-.18,-.22) rectangle (12.18,2.28);
\draw[black!42,line width=.35pt] (0,0) rectangle (12,2);
\foreach \x in {3,6,9}{\draw[black!38,dash pattern=on 1.2pt off 1.3pt,line width=.3pt] (\x,0)--(\x,2);}
\draw[redpen,line width=.9pt] (6,0)--(12,0)--(3,2)--(0,2)--cycle;
\draw[redpen,line width=.4pt,densely dashed] (1.5,0)--(6,0);
\draw[redpen,dash pattern=on 2pt off 1.6pt,line width=.55pt] (1.5,2)--(1.5,0);
\draw[redpen,line width=.4pt] (1.74,0)--(1.74,.24)--(1.5,.24);
\node[text=redpen,font=\ansfont\fontsize{8.6pt}{10pt}\selectfont,fill=white,inner sep=.5pt,anchor=west] at (1.98,1) {高};
\foreach \x in {0,3,6,9,12}{\foreach \y in {0,2}{\fill[black!70] (\x,\y) circle[radius=.075];}}
\end{tikzpicture}\par\vspace{.5mm}
上底\(3\)厘米，下底\(6\)厘米，高\(2\)厘米。\par 上边左端比下边左端向左错开\(6\)厘米。\end{minipage}
}{\begin{minipage}[t][44mm][t]{90mm}\ans\fontsize{10.6pt}{14pt}\selectfont
{\headfont\fontsize{11pt}{14pt}\selectfont 梯形\(2\)}\par\vspace{1mm}
\begin{tikzpicture}[x=6.3mm,y=6.3mm,line join=round,line cap=round]
\path[use as bounding box] (-.18,-.22) rectangle (12.18,2.28);
\draw[black!42,line width=.35pt] (0,0) rectangle (12,2);
\foreach \x in {3,6,9}{\draw[black!38,dash pattern=on 1.2pt off 1.3pt,line width=.3pt] (\x,0)--(\x,2);}
\draw[redpen,line width=.9pt] (3,0)--(9,0)--(3,2)--(0,2)--cycle;
\draw[redpen,line width=.4pt,densely dashed] (1.5,0)--(3,0);
\draw[redpen,dash pattern=on 2pt off 1.6pt,line width=.55pt] (1.5,2)--(1.5,0);
\draw[redpen,line width=.4pt] (1.74,0)--(1.74,.24)--(1.5,.24);
\node[text=redpen,font=\ansfont\fontsize{8.6pt}{10pt}\selectfont,fill=white,inner sep=.5pt,anchor=west] at (1.98,1) {高};
\foreach \x in {0,3,6,9,12}{\foreach \y in {0,2}{\fill[black!70] (\x,\y) circle[radius=.075];}}
\end{tikzpicture}\par\vspace{.5mm}
上底\(3\)厘米，下底\(6\)厘米，高\(2\)厘米。\par 上边左端比下边左端向左错开\(3\)厘米。\end{minipage}
}
\pair{\begin{minipage}[t][44mm][t]{90mm}\ans\fontsize{10.6pt}{14pt}\selectfont
{\headfont\fontsize{11pt}{14pt}\selectfont 梯形\(3\)}\par\vspace{1mm}
\begin{tikzpicture}[x=6.3mm,y=6.3mm,line join=round,line cap=round]
\path[use as bounding box] (-.18,-.22) rectangle (12.18,2.28);
\draw[black!42,line width=.35pt] (0,0) rectangle (12,2);
\foreach \x in {3,6,9}{\draw[black!38,dash pattern=on 1.2pt off 1.3pt,line width=.3pt] (\x,0)--(\x,2);}
\draw[redpen,line width=.9pt] (0,0)--(6,0)--(3,2)--(0,2)--cycle;
% Current revision: mark both right-angle vertices on the vertical left leg.
\draw[redpen,line width=.6bp] (.2,0)--(.2,.2)--(0,.2);
\draw[redpen,line width=.6bp] (.2,2)--(.2,1.8)--(0,1.8);

\draw[redpen,dash pattern=on 2pt off 1.6pt,line width=.55pt] (1.5,2)--(1.5,0);
\draw[redpen,line width=.4pt] (1.74,0)--(1.74,.24)--(1.5,.24);
\node[text=redpen,font=\ansfont\fontsize{8.6pt}{10pt}\selectfont,fill=white,inner sep=.5pt,anchor=west] at (1.98,1) {高};
\foreach \x in {0,3,6,9,12}{\foreach \y in {0,2}{\fill[black!70] (\x,\y) circle[radius=.075];}}
\end{tikzpicture}\par\vspace{.5mm}
上底\(3\)厘米，下底\(6\)厘米，高\(2\)厘米。\par 上边左端与下边左端对齐。\end{minipage}
}{\begin{minipage}[t][44mm][t]{90mm}\ans\fontsize{10.6pt}{14pt}\selectfont
{\headfont\fontsize{11pt}{14pt}\selectfont 梯形\(4\)}\par\vspace{1mm}
\begin{tikzpicture}[x=6.3mm,y=6.3mm,line join=round,line cap=round]
\path[use as bounding box] (-.18,-.22) rectangle (12.18,2.28);
\draw[black!42,line width=.35pt] (0,0) rectangle (12,2);
\foreach \x in {3,6,9}{\draw[black!38,dash pattern=on 1.2pt off 1.3pt,line width=.3pt] (\x,0)--(\x,2);}
\draw[redpen,line width=.9pt] (3,0)--(12,0)--(3,2)--(0,2)--cycle;
\draw[redpen,line width=.4pt,densely dashed] (1.5,0)--(3,0);
\draw[redpen,dash pattern=on 2pt off 1.6pt,line width=.55pt] (1.5,2)--(1.5,0);
\draw[redpen,line width=.4pt] (1.74,0)--(1.74,.24)--(1.5,.24);
\node[text=redpen,font=\ansfont\fontsize{8.6pt}{10pt}\selectfont,fill=white,inner sep=.5pt,anchor=west] at (1.98,1) {高};
\foreach \x in {0,3,6,9,12}{\foreach \y in {0,2}{\fill[black!70] (\x,\y) circle[radius=.075];}}
\end{tikzpicture}\par\vspace{.5mm}
上底\(3\)厘米，下底\(9\)厘米，高\(2\)厘米。\par 上边左端比下边左端向左错开\(3\)厘米。\end{minipage}
}
\pair{\begin{minipage}[t][44mm][t]{90mm}\ans\fontsize{10.6pt}{14pt}\selectfont
{\headfont\fontsize{11pt}{14pt}\selectfont 梯形\(5\)}\par\vspace{1mm}
\begin{tikzpicture}[x=6.3mm,y=6.3mm,line join=round,line cap=round]
\path[use as bounding box] (-.18,-.22) rectangle (12.18,2.28);
\draw[black!42,line width=.35pt] (0,0) rectangle (12,2);
\foreach \x in {3,6,9}{\draw[black!38,dash pattern=on 1.2pt off 1.3pt,line width=.3pt] (\x,0)--(\x,2);}
\draw[redpen,line width=.9pt] (0,0)--(9,0)--(3,2)--(0,2)--cycle;
% Current revision: mark both right-angle vertices on the vertical left leg.
\draw[redpen,line width=.6bp] (.2,0)--(.2,.2)--(0,.2);
\draw[redpen,line width=.6bp] (.2,2)--(.2,1.8)--(0,1.8);

\draw[redpen,dash pattern=on 2pt off 1.6pt,line width=.55pt] (1.5,2)--(1.5,0);
\draw[redpen,line width=.4pt] (1.74,0)--(1.74,.24)--(1.5,.24);
\node[text=redpen,font=\ansfont\fontsize{8.6pt}{10pt}\selectfont,fill=white,inner sep=.5pt,anchor=west] at (1.98,1) {高};
\foreach \x in {0,3,6,9,12}{\foreach \y in {0,2}{\fill[black!70] (\x,\y) circle[radius=.075];}}
\end{tikzpicture}\par\vspace{.5mm}
上底\(3\)厘米，下底\(9\)厘米，高\(2\)厘米。\par 上边左端与下边左端对齐。\end{minipage}
}{\begin{minipage}[t][44mm][t]{90mm}\ans\fontsize{10.6pt}{14pt}\selectfont
{\headfont\fontsize{11pt}{14pt}\selectfont 梯形\(6\)}\par\vspace{1mm}
\begin{tikzpicture}[x=6.3mm,y=6.3mm,line join=round,line cap=round]
\path[use as bounding box] (-.18,-.22) rectangle (12.18,2.28);
\draw[black!42,line width=.35pt] (0,0) rectangle (12,2);
\foreach \x in {3,6,9}{\draw[black!38,dash pattern=on 1.2pt off 1.3pt,line width=.3pt] (\x,0)--(\x,2);}
\draw[redpen,line width=.9pt] (0,0)--(9,0)--(6,2)--(3,2)--cycle;

\draw[redpen,dash pattern=on 2pt off 1.6pt,line width=.55pt] (4.5,2)--(4.5,0);
\draw[redpen,line width=.4pt] (4.74,0)--(4.74,.24)--(4.5,.24);
\node[text=redpen,font=\ansfont\fontsize{8.6pt}{10pt}\selectfont,fill=white,inner sep=.5pt,anchor=west] at (4.98,1) {高};
\foreach \x in {0,3,6,9,12}{\foreach \y in {0,2}{\fill[black!70] (\x,\y) circle[radius=.075];}}
\end{tikzpicture}\par\vspace{.5mm}
上底\(3\)厘米，下底\(9\)厘米，高\(2\)厘米。\par 上边左端比下边左端向右错开\(3\)厘米。\end{minipage}
}
\vspace{2mm}
\begin{minipage}[t]{89mm}\ans\fontsize{11pt}{16pt}\selectfont
\textbf{本页分类}\par
上底\(3\)厘米、下底\(6\)厘米：\(3\)种。\par
上底\(3\)厘米、下底\(9\)厘米：\(3\)种。\par
本页共\(3+3=6\)种，高均为\(2\)厘米。
\end{minipage}\hfill
\begin{minipage}[t]{89mm}\ans\fontsize{11pt}{16pt}\selectfont
\textbf{观察提醒}\par
两个梯形的上底、下底和高分别相等，\par
形状仍可能不同。\par
因此，不能只按上、下底的长度来数。\par
图中的高垂直于下底或下底的延长线。
\end{minipage}
\footer{1}
\newpage
\pagetitle{一、梯形（续：第\(7\)至第\(13\)种）}{2}
{\fontsize{10.5pt}{15pt}\selectfont
下面继续按较短底的长度分类。所有图形仍放在同样的\(12\)厘米长、\(2\)厘米宽的纸条中。\par}
\vspace{3mm}
\pair{\begin{minipage}[t][40mm][t]{90mm}\ans\fontsize{10.6pt}{14pt}\selectfont
{\headfont\fontsize{11pt}{14pt}\selectfont 梯形\(7\)}\par\vspace{1mm}
\begin{tikzpicture}[x=6.3mm,y=6.3mm,line join=round,line cap=round]
\path[use as bounding box] (-.18,-.22) rectangle (12.18,2.28);
\draw[black!42,line width=.35pt] (0,0) rectangle (12,2);
\foreach \x in {3,6,9}{\draw[black!38,dash pattern=on 1.2pt off 1.3pt,line width=.3pt] (\x,0)--(\x,2);}
\draw[redpen,line width=.9pt] (0,0)--(12,0)--(3,2)--(0,2)--cycle;
% Current revision: mark both right-angle vertices on the vertical left leg.
\draw[redpen,line width=.6bp] (.2,0)--(.2,.2)--(0,.2);
\draw[redpen,line width=.6bp] (.2,2)--(.2,1.8)--(0,1.8);

\draw[redpen,dash pattern=on 2pt off 1.6pt,line width=.55pt] (1.5,2)--(1.5,0);
\draw[redpen,line width=.4pt] (1.74,0)--(1.74,.24)--(1.5,.24);
\node[text=redpen,font=\ansfont\fontsize{8.6pt}{10pt}\selectfont,fill=white,inner sep=.5pt,anchor=west] at (1.98,1) {高};
\foreach \x in {0,3,6,9,12}{\foreach \y in {0,2}{\fill[black!70] (\x,\y) circle[radius=.075];}}
\end{tikzpicture}\par\vspace{.5mm}
上底\(3\)厘米，下底\(12\)厘米，高\(2\)厘米。\par 上边左端与下边左端对齐。\end{minipage}
}{\begin{minipage}[t][40mm][t]{90mm}\ans\fontsize{10.6pt}{14pt}\selectfont
{\headfont\fontsize{11pt}{14pt}\selectfont 梯形\(8\)}\par\vspace{1mm}
\begin{tikzpicture}[x=6.3mm,y=6.3mm,line join=round,line cap=round]
\path[use as bounding box] (-.18,-.22) rectangle (12.18,2.28);
\draw[black!42,line width=.35pt] (0,0) rectangle (12,2);
\foreach \x in {3,6,9}{\draw[black!38,dash pattern=on 1.2pt off 1.3pt,line width=.3pt] (\x,0)--(\x,2);}
\draw[redpen,line width=.9pt] (0,0)--(12,0)--(6,2)--(3,2)--cycle;

\draw[redpen,dash pattern=on 2pt off 1.6pt,line width=.55pt] (4.5,2)--(4.5,0);
\draw[redpen,line width=.4pt] (4.74,0)--(4.74,.24)--(4.5,.24);
\node[text=redpen,font=\ansfont\fontsize{8.6pt}{10pt}\selectfont,fill=white,inner sep=.5pt,anchor=west] at (4.98,1) {高};
\foreach \x in {0,3,6,9,12}{\foreach \y in {0,2}{\fill[black!70] (\x,\y) circle[radius=.075];}}
\end{tikzpicture}\par\vspace{.5mm}
上底\(3\)厘米，下底\(12\)厘米，高\(2\)厘米。\par 上边左端比下边左端向右错开\(3\)厘米。\end{minipage}
}
\pair{\begin{minipage}[t][40mm][t]{90mm}\ans\fontsize{10.6pt}{14pt}\selectfont
{\headfont\fontsize{11pt}{14pt}\selectfont 梯形\(9\)}\par\vspace{1mm}
\begin{tikzpicture}[x=6.3mm,y=6.3mm,line join=round,line cap=round]
\path[use as bounding box] (-.18,-.22) rectangle (12.18,2.28);
\draw[black!42,line width=.35pt] (0,0) rectangle (12,2);
\foreach \x in {3,6,9}{\draw[black!38,dash pattern=on 1.2pt off 1.3pt,line width=.3pt] (\x,0)--(\x,2);}
\draw[redpen,line width=.9pt] (3,0)--(12,0)--(6,2)--(0,2)--cycle;

\draw[redpen,dash pattern=on 2pt off 1.6pt,line width=.55pt] (3.0,2)--(3.0,0);
\draw[redpen,line width=.4pt] (3.24,0)--(3.24,.24)--(3.0,.24);
\node[text=redpen,font=\ansfont\fontsize{8.6pt}{10pt}\selectfont,fill=white,inner sep=.5pt,anchor=west] at (3.48,1) {高};
\foreach \x in {0,3,6,9,12}{\foreach \y in {0,2}{\fill[black!70] (\x,\y) circle[radius=.075];}}
\end{tikzpicture}\par\vspace{.5mm}
上底\(6\)厘米，下底\(9\)厘米，高\(2\)厘米。\par 上边左端比下边左端向左错开\(3\)厘米。\end{minipage}
}{\begin{minipage}[t][40mm][t]{90mm}\ans\fontsize{10.6pt}{14pt}\selectfont
{\headfont\fontsize{11pt}{14pt}\selectfont 梯形\(10\)}\par\vspace{1mm}
\begin{tikzpicture}[x=6.3mm,y=6.3mm,line join=round,line cap=round]
\path[use as bounding box] (-.18,-.22) rectangle (12.18,2.28);
\draw[black!42,line width=.35pt] (0,0) rectangle (12,2);
\foreach \x in {3,6,9}{\draw[black!38,dash pattern=on 1.2pt off 1.3pt,line width=.3pt] (\x,0)--(\x,2);}
\draw[redpen,line width=.9pt] (0,0)--(9,0)--(6,2)--(0,2)--cycle;
% Current revision: mark both right-angle vertices on the vertical left leg.
\draw[redpen,line width=.6bp] (.2,0)--(.2,.2)--(0,.2);
\draw[redpen,line width=.6bp] (.2,2)--(.2,1.8)--(0,1.8);

\draw[redpen,dash pattern=on 2pt off 1.6pt,line width=.55pt] (3.0,2)--(3.0,0);
\draw[redpen,line width=.4pt] (3.24,0)--(3.24,.24)--(3.0,.24);
\node[text=redpen,font=\ansfont\fontsize{8.6pt}{10pt}\selectfont,fill=white,inner sep=.5pt,anchor=west] at (3.48,1) {高};
\foreach \x in {0,3,6,9,12}{\foreach \y in {0,2}{\fill[black!70] (\x,\y) circle[radius=.075];}}
\end{tikzpicture}\par\vspace{.5mm}
上底\(6\)厘米，下底\(9\)厘米，高\(2\)厘米。\par 上边左端与下边左端对齐。\end{minipage}
}
\pair{\begin{minipage}[t][40mm][t]{90mm}\ans\fontsize{10.6pt}{14pt}\selectfont
{\headfont\fontsize{11pt}{14pt}\selectfont 梯形\(11\)}\par\vspace{1mm}
\begin{tikzpicture}[x=6.3mm,y=6.3mm,line join=round,line cap=round]
\path[use as bounding box] (-.18,-.22) rectangle (12.18,2.28);
\draw[black!42,line width=.35pt] (0,0) rectangle (12,2);
\foreach \x in {3,6,9}{\draw[black!38,dash pattern=on 1.2pt off 1.3pt,line width=.3pt] (\x,0)--(\x,2);}
\draw[redpen,line width=.9pt] (0,0)--(12,0)--(6,2)--(0,2)--cycle;
% Current revision: mark both right-angle vertices on the vertical left leg.
\draw[redpen,line width=.6bp] (.2,0)--(.2,.2)--(0,.2);
\draw[redpen,line width=.6bp] (.2,2)--(.2,1.8)--(0,1.8);

\draw[redpen,dash pattern=on 2pt off 1.6pt,line width=.55pt] (3.0,2)--(3.0,0);
\draw[redpen,line width=.4pt] (3.24,0)--(3.24,.24)--(3.0,.24);
\node[text=redpen,font=\ansfont\fontsize{8.6pt}{10pt}\selectfont,fill=white,inner sep=.5pt,anchor=west] at (3.48,1) {高};
\foreach \x in {0,3,6,9,12}{\foreach \y in {0,2}{\fill[black!70] (\x,\y) circle[radius=.075];}}
\end{tikzpicture}\par\vspace{.5mm}
上底\(6\)厘米，下底\(12\)厘米，高\(2\)厘米。\par 上边左端与下边左端对齐。\end{minipage}
}{\begin{minipage}[t][40mm][t]{90mm}\ans\fontsize{10.6pt}{14pt}\selectfont
{\headfont\fontsize{11pt}{14pt}\selectfont 梯形\(12\)}\par\vspace{1mm}
\begin{tikzpicture}[x=6.3mm,y=6.3mm,line join=round,line cap=round]
\path[use as bounding box] (-.18,-.22) rectangle (12.18,2.28);
\draw[black!42,line width=.35pt] (0,0) rectangle (12,2);
\foreach \x in {3,6,9}{\draw[black!38,dash pattern=on 1.2pt off 1.3pt,line width=.3pt] (\x,0)--(\x,2);}
\draw[redpen,line width=.9pt] (0,0)--(12,0)--(9,2)--(3,2)--cycle;

\draw[redpen,dash pattern=on 2pt off 1.6pt,line width=.55pt] (6.0,2)--(6.0,0);
\draw[redpen,line width=.4pt] (6.24,0)--(6.24,.24)--(6.0,.24);
\node[text=redpen,font=\ansfont\fontsize{8.6pt}{10pt}\selectfont,fill=white,inner sep=.5pt,anchor=west] at (6.48,1) {高};
\foreach \x in {0,3,6,9,12}{\foreach \y in {0,2}{\fill[black!70] (\x,\y) circle[radius=.075];}}
\end{tikzpicture}\par\vspace{.5mm}
上底\(6\)厘米，下底\(12\)厘米，高\(2\)厘米。\par 上边左端比下边左端向右错开\(3\)厘米。\end{minipage}
}
\pair{\begin{minipage}[t][40mm][t]{90mm}\ans\fontsize{10.6pt}{14pt}\selectfont
{\headfont\fontsize{11pt}{14pt}\selectfont 梯形\(13\)}\par\vspace{1mm}
\begin{tikzpicture}[x=6.3mm,y=6.3mm,line join=round,line cap=round]
\path[use as bounding box] (-.18,-.22) rectangle (12.18,2.28);
\draw[black!42,line width=.35pt] (0,0) rectangle (12,2);
\foreach \x in {3,6,9}{\draw[black!38,dash pattern=on 1.2pt off 1.3pt,line width=.3pt] (\x,0)--(\x,2);}
\draw[redpen,line width=.9pt] (0,0)--(12,0)--(9,2)--(0,2)--cycle;
% Current revision: mark both right-angle vertices on the vertical left leg.
\draw[redpen,line width=.6bp] (.2,0)--(.2,.2)--(0,.2);
\draw[redpen,line width=.6bp] (.2,2)--(.2,1.8)--(0,1.8);

\draw[redpen,dash pattern=on 2pt off 1.6pt,line width=.55pt] (4.5,2)--(4.5,0);
\draw[redpen,line width=.4pt] (4.74,0)--(4.74,.24)--(4.5,.24);
\node[text=redpen,font=\ansfont\fontsize{8.6pt}{10pt}\selectfont,fill=white,inner sep=.5pt,anchor=west] at (4.98,1) {高};
\foreach \x in {0,3,6,9,12}{\foreach \y in {0,2}{\fill[black!70] (\x,\y) circle[radius=.075];}}
\end{tikzpicture}\par\vspace{.5mm}
上底\(9\)厘米，下底\(12\)厘米，高\(2\)厘米。\par 上边左端与下边左端对齐。\end{minipage}
}{\begin{minipage}[t]{90mm}\mbox{}\end{minipage}}
\vspace{1mm}
\begin{minipage}[t]{89mm}\ans\fontsize{10.6pt}{15pt}\selectfont
\textbf{全部梯形分类汇总}\par
上底\(3\)厘米：下底\(6\)、\(9\)、\(12\)厘米，\par
分别有\(3\)、\(3\)、\(2\)种，共\(8\)种。\par
上底\(6\)厘米：下底\(9\)、\(12\)厘米，\par
分别有\(2\)、\(2\)种，共\(4\)种。\par
上底\(9\)厘米：下底\(12\)厘米，共\(1\)种。\par
合计：\(8+4+1=13\)种。
\end{minipage}\hfill
\begin{minipage}[t]{89mm}\ans\fontsize{10.6pt}{15pt}\selectfont
\textbf{为什么这些就全了}\par
底长只能是\(3\)、\(6\)、\(9\)、\(12\)厘米。\par
梯形的两条底不一样长，先把底长配全；\par
再让上底逐个移到允许的交点位置，\par
去掉平移、旋转或翻转后完全重合的图形，\par
就得到这\(13\)种。每种的高都是\(2\)厘米。
\end{minipage}
\footer{2}
\newpage
\pagetitle{二、平行四边形（全部\(10\)种，其中\(4\)种是长方形）}{3}
{\fontsize{10.3pt}{15pt}\selectfont
上下两条底取相同长度，再改变上下两条底的相对位置。\textbf{长方形也是平行四边形。}\par
下图所写的“底”都指水平边，所写的“高”都是这条底对应的垂直距离。\par}
\vspace{3mm}
\pair{\begin{minipage}[t][36mm][t]{90mm}\ans\fontsize{10.6pt}{14pt}\selectfont
{\headfont\fontsize{11pt}{14pt}\selectfont 平行四边形\(1\)}\quad（长方形）\par\vspace{1mm}
\begin{tikzpicture}[x=6.3mm,y=6.3mm,line join=round,line cap=round]
\path[use as bounding box] (-.18,-.22) rectangle (12.18,2.28);
\draw[black!42,line width=.35pt] (0,0) rectangle (12,2);
\foreach \x in {3,6,9}{\draw[black!38,dash pattern=on 1.2pt off 1.3pt,line width=.3pt] (\x,0)--(\x,2);}
\draw[redpen,line width=.9pt] (0,0)--(3,0)--(3,2)--(0,2)--cycle;

\draw[redpen,dash pattern=on 2pt off 1.6pt,line width=.55pt] (1.5,2)--(1.5,0);
\draw[redpen,line width=.4pt] (1.74,0)--(1.74,.24)--(1.5,.24);
\node[text=redpen,font=\ansfont\fontsize{8.6pt}{10pt}\selectfont,fill=white,inner sep=.5pt,anchor=west] at (1.98,1) {高};
\foreach \x in {0,3,6,9,12}{\foreach \y in {0,2}{\fill[black!70] (\x,\y) circle[radius=.075];}}
\end{tikzpicture}\par\vspace{.5mm}
底\(3\)厘米，高\(2\)厘米。\par 上边左端与下边左端对齐。\end{minipage}
}{\begin{minipage}[t][36mm][t]{90mm}\ans\fontsize{10.6pt}{14pt}\selectfont
{\headfont\fontsize{11pt}{14pt}\selectfont 平行四边形\(2\)}\par\vspace{1mm}
\begin{tikzpicture}[x=6.3mm,y=6.3mm,line join=round,line cap=round]
\path[use as bounding box] (-.18,-.22) rectangle (12.18,2.28);
\draw[black!42,line width=.35pt] (0,0) rectangle (12,2);
\foreach \x in {3,6,9}{\draw[black!38,dash pattern=on 1.2pt off 1.3pt,line width=.3pt] (\x,0)--(\x,2);}
\draw[redpen,line width=.9pt] (0,0)--(3,0)--(6,2)--(3,2)--cycle;
\draw[redpen,line width=.4pt,densely dashed] (3,0)--(4.5,0);
\draw[redpen,dash pattern=on 2pt off 1.6pt,line width=.55pt] (4.5,2)--(4.5,0);
\draw[redpen,line width=.4pt] (4.26,0)--(4.26,.24)--(4.5,.24);
\node[text=redpen,font=\ansfont\fontsize{8.6pt}{10pt}\selectfont,fill=white,inner sep=.5pt,anchor=west] at (4.98,1) {高};
\foreach \x in {0,3,6,9,12}{\foreach \y in {0,2}{\fill[black!70] (\x,\y) circle[radius=.075];}}
\end{tikzpicture}\par\vspace{.5mm}
底\(3\)厘米，高\(2\)厘米。\par 上边左端比下边左端向右错开\(3\)厘米。\end{minipage}
}
\pair{\begin{minipage}[t][36mm][t]{90mm}\ans\fontsize{10.6pt}{14pt}\selectfont
{\headfont\fontsize{11pt}{14pt}\selectfont 平行四边形\(3\)}\par\vspace{1mm}
\begin{tikzpicture}[x=6.3mm,y=6.3mm,line join=round,line cap=round]
\path[use as bounding box] (-.18,-.22) rectangle (12.18,2.28);
\draw[black!42,line width=.35pt] (0,0) rectangle (12,2);
\foreach \x in {3,6,9}{\draw[black!38,dash pattern=on 1.2pt off 1.3pt,line width=.3pt] (\x,0)--(\x,2);}
\draw[redpen,line width=.9pt] (0,0)--(3,0)--(9,2)--(6,2)--cycle;
\draw[redpen,line width=.4pt,densely dashed] (3,0)--(7.5,0);
\draw[redpen,dash pattern=on 2pt off 1.6pt,line width=.55pt] (7.5,2)--(7.5,0);
\draw[redpen,line width=.4pt] (7.26,0)--(7.26,.24)--(7.5,.24);
\node[text=redpen,font=\ansfont\fontsize{8.6pt}{10pt}\selectfont,fill=white,inner sep=.5pt,anchor=west] at (7.98,1) {高};
\foreach \x in {0,3,6,9,12}{\foreach \y in {0,2}{\fill[black!70] (\x,\y) circle[radius=.075];}}
\end{tikzpicture}\par\vspace{.5mm}
底\(3\)厘米，高\(2\)厘米。\par 上边左端比下边左端向右错开\(6\)厘米。\end{minipage}
}{\begin{minipage}[t][36mm][t]{90mm}\ans\fontsize{10.6pt}{14pt}\selectfont
{\headfont\fontsize{11pt}{14pt}\selectfont 平行四边形\(4\)}\par\vspace{1mm}
\begin{tikzpicture}[x=6.3mm,y=6.3mm,line join=round,line cap=round]
\path[use as bounding box] (-.18,-.22) rectangle (12.18,2.28);
\draw[black!42,line width=.35pt] (0,0) rectangle (12,2);
\foreach \x in {3,6,9}{\draw[black!38,dash pattern=on 1.2pt off 1.3pt,line width=.3pt] (\x,0)--(\x,2);}
\draw[redpen,line width=.9pt] (0,0)--(3,0)--(12,2)--(9,2)--cycle;
\draw[redpen,line width=.4pt,densely dashed] (3,0)--(10.5,0);
\draw[redpen,dash pattern=on 2pt off 1.6pt,line width=.55pt] (10.5,2)--(10.5,0);
\draw[redpen,line width=.4pt] (10.26,0)--(10.26,.24)--(10.5,.24);
\node[text=redpen,font=\ansfont\fontsize{8.6pt}{10pt}\selectfont,fill=white,inner sep=.5pt,anchor=west] at (10.98,1) {高};
\foreach \x in {0,3,6,9,12}{\foreach \y in {0,2}{\fill[black!70] (\x,\y) circle[radius=.075];}}
\end{tikzpicture}\par\vspace{.5mm}
底\(3\)厘米，高\(2\)厘米。\par 上边左端比下边左端向右错开\(9\)厘米。\end{minipage}
}
\pair{\begin{minipage}[t][36mm][t]{90mm}\ans\fontsize{10.6pt}{14pt}\selectfont
{\headfont\fontsize{11pt}{14pt}\selectfont 平行四边形\(5\)}\quad（长方形）\par\vspace{1mm}
\begin{tikzpicture}[x=6.3mm,y=6.3mm,line join=round,line cap=round]
\path[use as bounding box] (-.18,-.22) rectangle (12.18,2.28);
\draw[black!42,line width=.35pt] (0,0) rectangle (12,2);
\foreach \x in {3,6,9}{\draw[black!38,dash pattern=on 1.2pt off 1.3pt,line width=.3pt] (\x,0)--(\x,2);}
\draw[redpen,line width=.9pt] (0,0)--(6,0)--(6,2)--(0,2)--cycle;

\draw[redpen,dash pattern=on 2pt off 1.6pt,line width=.55pt] (3.0,2)--(3.0,0);
\draw[redpen,line width=.4pt] (3.24,0)--(3.24,.24)--(3.0,.24);
\node[text=redpen,font=\ansfont\fontsize{8.6pt}{10pt}\selectfont,fill=white,inner sep=.5pt,anchor=west] at (3.48,1) {高};
\foreach \x in {0,3,6,9,12}{\foreach \y in {0,2}{\fill[black!70] (\x,\y) circle[radius=.075];}}
\end{tikzpicture}\par\vspace{.5mm}
底\(6\)厘米，高\(2\)厘米。\par 上边左端与下边左端对齐。\end{minipage}
}{\begin{minipage}[t][36mm][t]{90mm}\ans\fontsize{10.6pt}{14pt}\selectfont
{\headfont\fontsize{11pt}{14pt}\selectfont 平行四边形\(6\)}\par\vspace{1mm}
\begin{tikzpicture}[x=6.3mm,y=6.3mm,line join=round,line cap=round]
\path[use as bounding box] (-.18,-.22) rectangle (12.18,2.28);
\draw[black!42,line width=.35pt] (0,0) rectangle (12,2);
\foreach \x in {3,6,9}{\draw[black!38,dash pattern=on 1.2pt off 1.3pt,line width=.3pt] (\x,0)--(\x,2);}
\draw[redpen,line width=.9pt] (0,0)--(6,0)--(9,2)--(3,2)--cycle;

\draw[redpen,dash pattern=on 2pt off 1.6pt,line width=.55pt] (6.0,2)--(6.0,0);
\draw[redpen,line width=.4pt] (5.76,0)--(5.76,.24)--(6.0,.24);
\node[text=redpen,font=\ansfont\fontsize{8.6pt}{10pt}\selectfont,fill=white,inner sep=.5pt,anchor=west] at (6.48,1) {高};
\foreach \x in {0,3,6,9,12}{\foreach \y in {0,2}{\fill[black!70] (\x,\y) circle[radius=.075];}}
\end{tikzpicture}\par\vspace{.5mm}
底\(6\)厘米，高\(2\)厘米。\par 上边左端比下边左端向右错开\(3\)厘米。\end{minipage}
}
\pair{\begin{minipage}[t][36mm][t]{90mm}\ans\fontsize{10.6pt}{14pt}\selectfont
{\headfont\fontsize{11pt}{14pt}\selectfont 平行四边形\(7\)}\par\vspace{1mm}
\begin{tikzpicture}[x=6.3mm,y=6.3mm,line join=round,line cap=round]
\path[use as bounding box] (-.18,-.22) rectangle (12.18,2.28);
\draw[black!42,line width=.35pt] (0,0) rectangle (12,2);
\foreach \x in {3,6,9}{\draw[black!38,dash pattern=on 1.2pt off 1.3pt,line width=.3pt] (\x,0)--(\x,2);}
\draw[redpen,line width=.9pt] (0,0)--(6,0)--(12,2)--(6,2)--cycle;
\draw[redpen,line width=.4pt,densely dashed] (6,0)--(9.0,0);
\draw[redpen,dash pattern=on 2pt off 1.6pt,line width=.55pt] (9.0,2)--(9.0,0);
\draw[redpen,line width=.4pt] (8.76,0)--(8.76,.24)--(9.0,.24);
\node[text=redpen,font=\ansfont\fontsize{8.6pt}{10pt}\selectfont,fill=white,inner sep=.5pt,anchor=west] at (9.48,1) {高};
\foreach \x in {0,3,6,9,12}{\foreach \y in {0,2}{\fill[black!70] (\x,\y) circle[radius=.075];}}
\end{tikzpicture}\par\vspace{.5mm}
底\(6\)厘米，高\(2\)厘米。\par 上边左端比下边左端向右错开\(6\)厘米。\end{minipage}
}{\begin{minipage}[t][36mm][t]{90mm}\ans\fontsize{10.6pt}{14pt}\selectfont
{\headfont\fontsize{11pt}{14pt}\selectfont 平行四边形\(8\)}\quad（长方形）\par\vspace{1mm}
\begin{tikzpicture}[x=6.3mm,y=6.3mm,line join=round,line cap=round]
\path[use as bounding box] (-.18,-.22) rectangle (12.18,2.28);
\draw[black!42,line width=.35pt] (0,0) rectangle (12,2);
\foreach \x in {3,6,9}{\draw[black!38,dash pattern=on 1.2pt off 1.3pt,line width=.3pt] (\x,0)--(\x,2);}
\draw[redpen,line width=.9pt] (0,0)--(9,0)--(9,2)--(0,2)--cycle;

\draw[redpen,dash pattern=on 2pt off 1.6pt,line width=.55pt] (4.5,2)--(4.5,0);
\draw[redpen,line width=.4pt] (4.74,0)--(4.74,.24)--(4.5,.24);
\node[text=redpen,font=\ansfont\fontsize{8.6pt}{10pt}\selectfont,fill=white,inner sep=.5pt,anchor=west] at (4.98,1) {高};
\foreach \x in {0,3,6,9,12}{\foreach \y in {0,2}{\fill[black!70] (\x,\y) circle[radius=.075];}}
\end{tikzpicture}\par\vspace{.5mm}
底\(9\)厘米，高\(2\)厘米。\par 上边左端与下边左端对齐。\end{minipage}
}
\pair{\begin{minipage}[t][36mm][t]{90mm}\ans\fontsize{10.6pt}{14pt}\selectfont
{\headfont\fontsize{11pt}{14pt}\selectfont 平行四边形\(9\)}\par\vspace{1mm}
\begin{tikzpicture}[x=6.3mm,y=6.3mm,line join=round,line cap=round]
\path[use as bounding box] (-.18,-.22) rectangle (12.18,2.28);
\draw[black!42,line width=.35pt] (0,0) rectangle (12,2);
\foreach \x in {3,6,9}{\draw[black!38,dash pattern=on 1.2pt off 1.3pt,line width=.3pt] (\x,0)--(\x,2);}
\draw[redpen,line width=.9pt] (0,0)--(9,0)--(12,2)--(3,2)--cycle;

\draw[redpen,dash pattern=on 2pt off 1.6pt,line width=.55pt] (7.5,2)--(7.5,0);
\draw[redpen,line width=.4pt] (7.74,0)--(7.74,.24)--(7.5,.24);
\node[text=redpen,font=\ansfont\fontsize{8.6pt}{10pt}\selectfont,fill=white,inner sep=.5pt,anchor=west] at (7.98,1) {高};
\foreach \x in {0,3,6,9,12}{\foreach \y in {0,2}{\fill[black!70] (\x,\y) circle[radius=.075];}}
\end{tikzpicture}\par\vspace{.5mm}
底\(9\)厘米，高\(2\)厘米。\par 上边左端比下边左端向右错开\(3\)厘米。\end{minipage}
}{\begin{minipage}[t][36mm][t]{90mm}\ans\fontsize{10.6pt}{14pt}\selectfont
{\headfont\fontsize{11pt}{14pt}\selectfont 平行四边形\(10\)}\quad（长方形）\par\vspace{1mm}
\begin{tikzpicture}[x=6.3mm,y=6.3mm,line join=round,line cap=round]
\path[use as bounding box] (-.18,-.22) rectangle (12.18,2.28);
\draw[black!42,line width=.35pt] (0,0) rectangle (12,2);
\foreach \x in {3,6,9}{\draw[black!38,dash pattern=on 1.2pt off 1.3pt,line width=.3pt] (\x,0)--(\x,2);}
\draw[redpen,line width=.9pt] (0,0)--(12,0)--(12,2)--(0,2)--cycle;

\draw[redpen,dash pattern=on 2pt off 1.6pt,line width=.55pt] (6.0,2)--(6.0,0);
\draw[redpen,line width=.4pt] (6.24,0)--(6.24,.24)--(6.0,.24);
\node[text=redpen,font=\ansfont\fontsize{8.6pt}{10pt}\selectfont,fill=white,inner sep=.5pt,anchor=west] at (6.48,1) {高};
\foreach \x in {0,3,6,9,12}{\foreach \y in {0,2}{\fill[black!70] (\x,\y) circle[radius=.075];}}
\end{tikzpicture}\par\vspace{.5mm}
底\(12\)厘米，高\(2\)厘米。\par 上边左端与下边左端对齐。\end{minipage}
}
\vspace{1mm}
\begin{minipage}[t]{89mm}\ans\fontsize{10.2pt}{14pt}\selectfont
\textbf{分类计数}\par
底\(3\)厘米有\(4\)种；底\(6\)厘米有\(3\)种；\par
底\(9\)厘米有\(2\)种；底\(12\)厘米有\(1\)种。\par
合计：\(4+3+2+1=10\)种，高均为\(2\)厘米。
\end{minipage}\hfill
\begin{minipage}[t]{89mm}\ans\fontsize{10.2pt}{14pt}\selectfont
\textbf{避免重复与遗漏}\par
相应的左移图与右移图翻转后完全重合，只算一种。\par
第\(1\)、\(5\)、\(8\)、\(10\)种是长方形。\par
若另外要求只画斜的，则有\(10-4=6\)种。
\end{minipage}
\footer{3}
\end{document}
